\chapter{Lean Metaprogramming: Generating the Closed-Form Proofs}
\label{ch:lean-meta}

\section{Overview}
\label{sec:meta-ov}

The Lean~4 layer (\cref{ch:lean}) does more than \emph{store} the closed
forms; it \emph{generates and checks} their proofs through Lean's
metaprogramming facilities. This chapter describes how the
\texttt{TheoremCalc} module uses tactics and the metaprogramming API to
turn an arithmetic identity into a certified theorem, and how the same
machinery reports open proof debt via \texttt{sorrydb}.

\section{The TheoremCalc Module}
\label{sec:meta-module}

\texttt{TheoremCalc} is a Lean package whose public surface is a small set
of theorems, each of the shape
\[
\forall n:\mathbb{N},\; \sum_{k=1}^{n} P(k) = C(n),
\]
where $P$ is a polynomial and $C$ its closed form. The module imports the
standard \texttt{Mathlib} algebraic hierarchies so that ring and
linear-arithmetic tactics apply.

\section{Proof by Reflection}
\label{sec:meta-reflect}

For polynomial summands, the proof is generated by \emph{reflection}: the
summand $P$ is represented as a concrete polynomial term, its summation is
computed symbolically in the metaprogram, and the result is reified back
into a Lean term that is then discharged by the \texttt{ring} tactic. The
pattern is:
\begin{enumerate}
  \item Parse the summand into a polynomial representation.
  \item Apply the known summation rules for monomials
        ($\sum k$, $\sum k^2$, $\sum k^3$, \dots).
  \item Assemble the closed form $C(n)$ in the meta-world.
  \item Emit a Lean proof term asserting
        $\sum P = C(n)$ and close it with \texttt{ring}.
\end{enumerate}

\section{Tactics Used}
\label{sec:meta-tactics}

The generated proofs rely on a small, well-audited tactic set:
\begin{itemize}
  \item \texttt{induction} on $n$ for the base-case / step structure.
  \item \texttt{ring} for the equational reasoning on polynomial
        expressions.
  \item \texttt{linear\_combination} to combine previously proven
        monomial identities into the target.
  \item \texttt{omega} / \texttt{linarith} for the index-range
        inequalities.
\end{itemize}
Because these tactics are themselves part of Lean's trusted core only
through their well-typed proof terms, the \emph{resulting} theorem carries
no more trust than a hand-written one, while the \emph{effort} of producing
it is mechanized.

\section{The \texttt{sorry} Debt and \texttt{sorrydb}}
\label{sec:meta-sorry}

When an identity cannot yet be generated (for example, a non-polynomial
summand), the module records a \texttt{sorry} placeholder and appends the
gap to \texttt{sorrydb}. A \texttt{sorry} is \emph{not} a silent pass: the
build gate (\cref{ch:build}) reads \texttt{sorrydb} and treats any entry
as an open item that lowers the completeness score $C$
(\cref{ch:completeness}). This converts "we haven't proven this yet" into a
countable, reviewable artifact rather than a hidden assumption.

\section{Connecting to the P/NP Swarm}
\label{sec:meta-swarm}

Each \texttt{sorrydb} entry becomes a candidate problem for the P/NP swarm
(\cref{ch:pnp}). An agent claims the problem, computes a witness (the
missing proof), and submits it; if Lean accepts the submitted proof term,
the \texttt{sorry} is replaced by a real theorem and $C$ increases. The
metaprogramming layer is therefore the bridge between human/agent proof
effort and the certified corpus: it both \emph{produces} proofs
automatically and \emph{consumes} externally supplied ones.

\section{Soundness Boundary of Metaprogramming}
\label{sec:meta-sound}

Lean's metaprogramming runs in a separate, unchecked monad; the
\emph{generated proof term}, however, is checked by the kernel before it
becomes a theorem. This separation is essential: bugs in the generator can
at worst fail to produce a proof (or produce one the kernel rejects), never
a false theorem. The trusted base is thus the Lean kernel plus the
reflection correctness, not the generator code itself.

\section{Example: the Cubic-Sum Identity}
\label{sec:meta-example}

For $P(k)=k^3$, the metaprogram assembles
\[
\sum_{k=1}^{n} k^3 = \left(\frac{n(n+1)}{2}\right)^2,
\]
proves it by induction with \texttt{ring} closing the inductive step, and
registers it as \texttt{TheoremCalc.cubeSum}. This is one of the three
closed forms currently in the certified corpus; the other two follow the
identical generated pattern for $P(k)=k$ and $P(k)=k^2$.
